%%%%%%%%%%%%%%%%%%%%%%%%%%%%%%%%%%%%%%%%%%%%%%%%%%%%%%%%%%%%%%%%%%%%%%%%%%%%%%%%%%%%%%%%%%%%%%%
%% Credit: Initial version written by [Ivan C. Christov](http://christov.tmnt-lab.org),      %%
%% Purdue University.                                                                        %%
%% License: GPLv3 (see https://www.gnu.org/licenses/gpl-3.0.en.html)                         %%
%%%%%%%%%%%%%%%%%%%%%%%%%%%%%%%%%%%%%%%%%%%%%%%%%%%%%%%%%%%%%%%%%%%%%%%%%%%%%%%%%%%%%%%%%%%%%%%
\documentclass[11pt]{article}
\usepackage[margin=0.75in,letterpaper]{geometry}
\usepackage{fancyhdr,totpages,parskip,enumitem}
%\setlist{nosep}

\usepackage{siunitx}
\usepackage{amsmath,bm}
\usepackage{fontawesome5}
\usepackage{tikzsymbols}

\usepackage[scaled]{helvet}
\usepackage[T1]{fontenc}
\renewcommand\familydefault{\sfdefault}

\pagestyle{fancyplain}
\fancyhf{}
%\renewcommand{\headrulewidth}{0pt}
\fancyhead[L]{\footnotesize Purdue University, ME 50900}
\fancyhead[R]{\footnotesize Prof.\ Ivan C.\ Christov}
\fancyfoot[R]{\footnotesize Page \thepage\ of \ref*{TotPages}}
\fancyfoot[L]{\footnotesize Last updated on \today\ by ICC}

\usepackage[colorlinks=true,urlcolor=blue]{hyperref}
\urlstyle{same}

\begin{document}

\textbf{Writing Tips and Best Practices for Assignments}

In science and engineering, writing is a primary means of communicating ideas. Whether it's project reports, journal articles, proposals, manuals, or even emails and memos, clear writing helps others understand your work. As a graduate student, you'll likely find that more people read your writing than hear you present. That means even the most innovative ideas can be overlooked if they aren't explained clearly and professionally. Strong technical writing skills are essential for effectively sharing your research, making a positive impression, and advancing in both academia and industry.

In this course, you'll be communicating your knowledge of the material in written form via written assignments, such as problem sets. These assignments present you with an ideal opportunity to practice your technical communication skills.

I encourage you to emulate the format used in the solutions to the problem sets that I will provide. These solutions have been carefully prepared, including typesetting the equations and providing graphical illustrations and references where needed to make them thorough and easy to follow. Thus, I hope they stand as good examples of high-quality technical writing.

\underline{Some tips and best practices:}

\begin{enumerate}

	\item Show your work, write neatly, and explain logically. Although \textbf{it is not required for this class}, it is recommended. However, typesetting reports is expected in professional practice. Try using \href{https://www.overleaf.com/edu/purdue}{Overleaf} (based on the powerful gold-standard \LaTeX\ typesetting language, which you can learn to use in \href{https://www.overleaf.com/learn/latex/Learn_LaTeX_in_30_minutes}{30 minutes}). It would be beneficial for you to learn how to use such software effectively.

	\item\label{itm:explain} Explain your solution process. Don't just list many equations on the page, and sprint towards a potential answer. Instead, \href{https://pubs.aip.org/asa/poma/article/54/1/025002/3352466/Helping-students-learn-to-read-mathematical}{explain what the equations mean in words}, then guide the reader (grader) through how you're manipulating the equations to arrive at the solution (often a symbolic expression in this course).

	\item Don't skip too many steps. You're more likely to make mistakes that way, making it harder for the reader (grader) to follow. Err on the side of more detail.

	\item Always start from the basic equations and simplify from there using assumptions and givens.

	\item Clearly identify significant assumptions you've made, and when and where you've used these assumptions (or givens) to make a simplification. Going back to item \ref{itm:explain}, the idea is to logically and unambiguously explain how you arrived at the answer.

	\item Consider using figures, diagrams, or schematics in your solution. Most of us absorb information more effectively through visual means. As needed, sketch the problem setup. Use high-quality plots to support your conclusions. Although it is helpful to have figures close to the text/results they are contributing, don't stress over precise location, especially if using \LaTeX, which has its own sophisticated rules about what placement ``looks best.''

	\item Define your variables. Be consistent in their use. Ideally, don't stray too far from the notation used in class or on the assignment.

	\item Include units with your numerical values. Verify that your unit conversions are correct.

	\item Clearly identify when a variable is a scalar (such as $f(x)$, $5$, or $b$), vector (such as $\underline{v}$, $\vec{v}$, or $\bm{v}$), or a tensor (such as $\underline{\underline{T}}$, $\vec{\vec{T}}$, or $\bm{\mathsf{T}}$). Just like with units, your use of scalars, vectors, and tensors in algebraic equations should be consistent.

	\item If typing your assignment:
	\begin{enumerate}
		\item Input mathematical symbols and expressions in ``math mode'' to ensure they are typeset correctly (usually that means in \textit{italics}). If using Word, use ``insert equation.'' If using \LaTeX, use dollar signs \verb+$...$+ or \verb+\begin{ } ... \end{ }+ environments to enclose the math expression. For example, ensure your output is ``The shear rate $\dot{\gamma}$ increases as ...'' instead of ``The shear rate y increases as ...'' In this example, I'm also illustrating the proper use of Greek letters; avoid using a ``close enough'' Latin letter. \Smiley

		\item If you want to learn about advanced equation typesetting, I recommend the \href{https://ctan.org/pkg/short-math-guide}{Short Math Guide for \LaTeX}.

		\item Don't italicize units. Units are ``text,'' they are not ``math.'' That is, write ``$101.3~\si{\kilo\pascal}$'' instead of ``$101.3~kPa$.'' Furthermore, ensure there is a space between numbers and units, that is, write ``$101.3~\si{\kilo\pascal}$'' instead of ``$101.3\si{\kilo\pascal}$.''

		\begin{enumerate}
			\item[\faLightbulb] Consider using the \LaTeX\ package \href{https://ctan.org/pkg/siunitx}{\texttt{siunitx}}, which allows you to typeset units as \verb+\si{\name}+. That is, write ``\verb+$101.3~\si{\kilo\pascal}$+'' instead of ``\verb+$101.3~kPa$+'' or ``\verb+$101.3$ kPa+.''

			\item[\faLightbulb] In \LaTeX\ math mode, spaces don't matter. If you want to ensure two symbols are separated appropriately, use a tilde (\verb+~+) to connect them, that is, write ``\verb+$101.3~\si{\kilo\pascal}$+'' instead of  ``\verb+$101.3 \si{\kilo\pascal}$+.'' \\
			Furthermore, in \LaTeX, \verb+~+ signifies and \emph{non-breaking} space. You don't want your units to end up on the next line after the value. In Word, a non-breaking space is inserted via the key combination CTRL+SHIFT+SPACE.

    		\item[\faLightbulb] Note that in \LaTeX\ left quotes are \verb+`+ (or coded as \verb+\lq+), which is a different key than \verb+'+. To get double quotation marks, press \verb+`+ twice, \verb+``+. To close the quote, press \verb+'+ twice, \verb+''+.

			\item[\faLightbulb] Also note that \LaTeX\ assumes periods denote the end of a sentence, which requires extra space after the period. However, we often use periods in the middle of a sentence, in constructions such as ``Eq.'' or ``et al.'' To ensure proper spacing, code these expressions as ``\verb+Eq.\+'' or ``\verb+et al.\+'' the backslash after the period tells \LaTeX\ that this period is not the end of a sentence.
		\end{enumerate}

	\end{enumerate}

	\item Use equation numbers if you want to refer back to them later in your solution. For example, ``Substituting $x=3$ in Eq.~(5), we arrive at ...'' If using \LaTeX, give your equations labels using \verb+\label{name}+ and reference them using \verb+\eqref{name}+ (available from the \href{https://ctan.org/pkg/amsmath}{\texttt{amsmath}} package) for automatic numbering. If you are using Word, you're on your own...good luck. \Smiley

	\item Number pages in your document.

	\item As Dr.\ Nicole Sharp points out in her article ``Adopting a communication lifestyle,'' \textit{Physical Review Fluids} \textbf{5} (2020) 110515, \href{https://journals-aps-org.ezproxy.lib.purdue.edu/prfluids/abstract/10.1103/PhysRevFluids.5.110515}{doi:10.1103/PhysRevFluids.5.110515} (Purdue login required), craft your submission by
	\begin{enumerate}
		\item identifying big-picture issues around goals, audience, and message;
		\item constructing a first version of the product;
		\item evaluating and refining the product based on the bigger picture and outside feedback; and
		\item reflecting on the lessons learned during the project.
	\end{enumerate}
	In other words, review your final written solution, think about how it might look from the reader's (grader's) perspective, and revise as needed.

	\item Consider doing a basic grammar check of your writing using a tool such as LanguageTool, Writefull, or Grammarly, all of which have a \emph{free} version and integrate with Overleaf and/or via browser extensions.

	\item When preparing plots:

	\begin{enumerate}
		\item Make sure the size of axes, labels, legends, etc.\ is easy to read. Don't use bold fonts unless you're specifically denoting vectors or tensors.

		\item Make sure your plot symbols are clear and discernible. Don't reuse plot symbols for multiple datasets.

		\item Don't put a title on the plot itself; any discussion of the plot (or what it represents) should be placed below it as a figure caption.

		\item Do not take screenshots, as they can be of low quality and pixelated.  Use the plotting software's built-in function (or button) to save a plot in a high-quality vector format, such as PDF.

		\item Label your axes, plot symbols, and plot curves.

		\item Make sure the legend doesn't cover the plot.

		\item Use scientific notation instead of coding notation. For example, write ``$F=1.6\times10^{7}~\si{\newton}$'' instead of ``$F = 1.6e7~\si{\newton}$.'' Note that $\times$ (coded \verb+$\times$+) is a special symbol; it is \emph{not} the math variable $x$ (coded \verb+$x$+) in a different font.

		\item Generally, use solid lines (or curves) for theory/predicted values and symbols for discrete data obtained from experiments/simulations. To avoid confusion, don't connect data points from experiments/simulations with line segments.

		\item If your figure contains multiple subplots, say (a), (b), (c), etc., and make sure that the axis ranges are the same across subplots if you are comparing the same quantity across them.
	\end{enumerate}

	\item Tables should be reserved for cataloging many sets of values relevant to the problem, such as the thermophysical properties of the fluid, or the dimensions of the flow geometry. Generally, we will not need to do that in this course. Tables should be used sparingly. Output data and results are more effectively displayed as figures rather than tables.

	\item Never use Microsoft Excel charts (especially the 3D ones and related atrocities) to convey quantitative scientific data.

	\item Further reading:

	\begin{enumerate}
		\item ``\href{http://www.cs.ucdavis.edu/~amenta/w10/writingman.pdf}{A Guide to Writing Mathematics}'' by Dr.\ Kevin P.\ Lee.

		\item \href{https://nhigham.com/handbook-of-writing-for-the-mathematical-sciences/}{\textit{Handbook of Writing for the Mathematical Sciences}} by Prof.\ Nicholas J.\ Higham, available for download as PDF through \href{https://purdue.primo.exlibrisgroup.com/permalink/01PURDUE_PUWL/ufs51j/alma99169877700501081}{Purdue Libraries} (login required).

		\item The Purdue Online Writing Laboratory: \url{https://owl.purdue.edu/}
	\end{enumerate}

\end{enumerate}

\end{document}
